%%% FIELD selected-statement-proposed
Let \((Q_1,Q_2,A_1,A_2)\) be the centered Gaussian limit of the pre-period loss and ATT coordinates, and let \((R_1,R_2)\) be the centered Gaussian post-period limit, independent of the pre-period limit.  For each \(s\in\mathcal S\), the fitted-score Bartlett estimator with \(q_n=\lfloor\min(n,m)^{1/5}\rfloor\) satisfies \(\widehat V_s^{\mathrm{conv}}\to_{P_n}V_s\) and \(\Pr_{P_n}\{\widehat V_s^{\mathrm{conv}}>0\}\to1\).  Consequently,
\[
 \sqrt n(\widehat\tau_{\mathrm{sel}}-\tau)
 \Rightarrow A_{J(d)}+\sqrt rR_{J(d)},
\]
and
\[
 \frac{\widehat\tau_{\mathrm{sel}}-\tau}
 {\widehat{\operatorname{se}}^{\mathrm{conv}}_{\widehat J}}
 \Rightarrow
 T_1\mathbf 1\{X\leq k\}+T_2\mathbf 1\{X>k\}.
\]
On the explicit common stationary square subclass certified by \(\texttt{prop:shi-anchor-transfer}\), fixing a predeclared role, or considering a separate sequence for which \(|\sqrt n(L_1-L_2)|\to\infty\), reduces these displays to the corresponding fixed-role centered normal law.  This is a common-subclass reduction, not a set-inclusion claim for the full published class or a claim about its rectangular application procedure.
%%% FIELD selected-statement-replacement
Let \((Q_1,Q_2,A_1,A_2)\) be the centered Gaussian limit of the pre-period loss and ATT coordinates, and let \((R_1,R_2)\) be the centered Gaussian post-period limit, independent of the pre-period limit.  For each \(s\in\mathcal S\), the fitted-score Bartlett estimator with \(q_n=\lfloor\min(n,m)^{1/5}\rfloor\) satisfies \(\widehat V_s^{\mathrm{conv}}\to_{P_n}V_s\) and \(\Pr_{P_n}\{\widehat V_s^{\mathrm{conv}}>0\}\to1\).  Consequently,
\[
 \sqrt n(\widehat\tau_{\mathrm{sel}}-\tau)
 \Rightarrow A_{J(d)}+\sqrt rR_{J(d)},
\]
and
\[
 \frac{\widehat\tau_{\mathrm{sel}}-\tau}
 {\widehat{\operatorname{se}}^{\mathrm{conv}}_{\widehat J}}
 \Rightarrow
 T_1\mathbf 1\{X\leq k\}+T_2\mathbf 1\{X>k\}.
\]
On the explicit common stationary square subclass certified by \(\texttt{prop:shi-anchor-transfer}\), fixing a predeclared role, or considering a separate sequence for which \(|\sqrt n(L_1-L_2)|\to\infty\), reduces these displays to the corresponding fixed-role centered normal law.  This is a common-subclass reduction, not a set-inclusion claim for the full published class or a claim about its rectangular application procedure.
%%% FIELD ridge-expansion-statement
Jointly for \(s\in\mathcal S\),
\[
 \sqrt n(\widehat L_s-L_s)
 =n^{-1/2}\sum_{t=1}^n(\ell_{s,t}-L_s)+o_{P_n}(1).
\]
%%% FIELD ridge-expansion-proof
Fix \(s\in\mathcal S\).  Pre-treatment consistency gives \(Y_t=Y_t(0)\) for \(1\leq t\leq n\).  Define the empirical ridge moments
\[
 \widehat K_s=I_{d_s}+\frac1n\sum_{t=1}^n
 W_{D_s,t}W_{D_s,t}^{\prime},
 \qquad
 \widehat b_s=\frac1n\sum_{t=1}^nW_{D_s,t}Y_t(0),
\]
and the empirical criterion
\[
 \widehat Q_s(a)=\frac1n\sum_{t=1}^n
 \{Y_t(0)-W_{D_s,t}^{\prime}a\}^2+\lVert a\rVert_2^2.
\]
Because \(\widehat K_s\succeq I_{d_s}\), the criterion is strictly convex and its unique minimizer is
\[
 \widehat a_s=\widehat K_s^{-1}\widehat b_s.
\]
Thus the fallback in \(\texttt{def:ridge-selector}\) is never invoked for the ridge fit.

We first record the stochastic orders used below.  If \(U_t\) and \(V_t\) are centered scalar functions of observations separated by \(j\) periods and have uniformly bounded \((2+\epsilon)\)-moments for some \(\epsilon>0\), truncating both variables at a level \(R\), applying the defining total-variation bound for the absolute-regularity coefficient to the bounded parts, and applying H\"older's and Markov's inequalities to the tails gives
\[
 |\operatorname{Cov}_{P_n}(U_t,V_{t-j})|
 \leq C\,\beta_n(j)^{\epsilon/(2+\epsilon)}.
\]
The constant depends only on the common moment bound.  Under \(\texttt{ass:geometric-beta-mixing}\), the right-hand side is bounded by \(C_1e^{-c_1j}\), uniformly in \(n\).  The fixed dimension and \(\texttt{ass:moment-bound}\) therefore imply, entrywise and hence in Euclidean or operator norm,
\[
 \widehat K_s-K_s=O_{P_n}(n^{-1/2}),
 \qquad
 \widehat b_s-b_s=O_{P_n}(n^{-1/2}),
 \qquad
 \lVert\widehat K_s\rVert_{\mathrm{op}}=O_{P_n}(1).
\]
Indeed, the variance of each centered sample mean is at most
\[
 \frac1n\left\{\gamma(0)+2\sum_{j\geq1}|\gamma(j)|\right\}
 \leq \frac Cn,
\]
and Chebyshev's inequality applies.  Here the required moments exist because entries of \(W_{D_s,t}W_{D_s,t}^{\prime}\) and \(W_{D_s,t}Y_t(0)\) are quadratic functions of the primitive vector and \(8+\delta_0>4\).  Moreover, \(K_s\succeq I_{d_s}\), so \(\lVert K_s^{-1}\rVert_{\mathrm{op}}\leq1\), while Cauchy--Schwarz and the moment bound give \(\lVert b_s\rVert_2\leq C\).  Hence \(\lVert a_s\rVert_2\leq C\).

The empirical first-order condition and \(K_sa_s=b_s\) now give the exact relation
\[
 \widehat a_s-a_s
 =\widehat K_s^{-1}
 \{(\widehat b_s-b_s)-(\widehat K_s-K_s)a_s\},
\]
and therefore
\[
 \widehat a_s-a_s=O_{P_n}(n^{-1/2}).
\]
Completing the quadratic square about its empirical minimizer gives another exact identity,
\[
 \widehat L_s
 =\widehat Q_s(a_s)
  -(\widehat a_s-a_s)^{\prime}\widehat K_s
    (\widehat a_s-a_s).
\]
By the preceding bounds, the last term is \(O_{P_n}(n^{-1})\).  On the other hand, the definition of \(\ell_{s,t}\) gives
\[
 \widehat Q_s(a_s)-L_s
 =\frac1n\sum_{t=1}^n(\ell_{s,t}-L_s).
\]
Multiplication by \(\sqrt n\) proves
\[
 \sqrt n(\widehat L_s-L_s)
 =n^{-1/2}\sum_{t=1}^n(\ell_{s,t}-L_s)+o_{P_n}(1).
\]
There are exactly two fixed-dimensional roles.  Taking the maximum of the two remainder probabilities proves the asserted joint expansion.
%%% FIELD selected-proof
By \(\texttt{lem:ridge-loss-expansion}\), jointly for the two roles,
\[
 \sqrt n\{(\widehat L_1-\widehat L_2)-(L_1-L_2)\}
 =n^{-1/2}\sum_{t=1}^n
 \{(\ell_{1,t}-L_1)-(\ell_{2,t}-L_2)\}+o_{P_n}(1).
\]
The Gaussian-limit premise in the theorem, \(\texttt{ass:local-loss-gap}\), and the definition \(D=Q_1-Q_2\) consequently imply the joint convergence
\[
 \sqrt n(\widehat L_1-\widehat L_2)\Rightarrow d+D.
\]
Assumption \(\texttt{ass:selector-variance-margin}\) states \(\operatorname{Var}(D)\geq\nu>0\).  Thus \(d+D\) has a continuous normal distribution even when the full covariance matrix of all displayed Gaussian coordinates is singular, and
\[
 \Pr\{|d+D|\leq\epsilon\}
 \leq \frac{2\epsilon}{\sqrt{2\pi\nu}}
 \qquad(\epsilon>0).
\]
Approximating the indicator of \(( -\infty,0]\) from above and below by continuous functions that differ only on \([-epsilon,epsilon]\), and then sending \(\epsilon\downarrow0\), proves from the tie rule in \(\texttt{def:ridge-selector}\) that
\[
 \widehat J\Rightarrow J(d),
 \qquad
 J(d)=1\ \text{on }\{d+D\leq0\},
 \quad
 J(d)=2\ \text{on }\{d+D>0\}.
\]

For each role, \(\texttt{lem:fixed-role-linearization}\) gives
\[
 \sqrt n(\widehat\tau_s-\tau)
 =n^{-1/2}\sum_{t=1}^n\psi_{s,t}
  +\sqrt{\frac nm}\,m^{-1/2}
    \sum_{t=n+1}^{n+m}v_{s,t}+o_{P_n}(1).
\]
The theorem's joint Gaussian premise and the explicitly required ratio condition \(n/m\to r\in(0,\infty)\) therefore yield
\[
 \left(
  \sqrt n(\widehat\tau_1-\tau),
  \sqrt n(\widehat\tau_2-\tau),
  \sqrt n(\widehat L_1-\widehat L_2)
 \right)
 \Rightarrow
 \left(A_1+\sqrt rR_1,A_2+\sqrt rR_2,d+D\right).
\]
The switching map is continuous except where its last coordinate is zero.  The preceding normal-density bound makes that discontinuity set null.  Equivalently, continuous upper and lower interpolations outside an \(\epsilon\)-strip, followed by tightness of the two ATT coordinates and then \(\epsilon\downarrow0\), give
\[
 \begin{aligned}
 \sqrt n(\widehat\tau_{\mathrm{sel}}-\tau)
 &=\sum_{s=1}^2\mathbf 1\{\widehat J=s\}
   \sqrt n(\widehat\tau_s-\tau)\\
 &\Rightarrow A_{J(d)}+\sqrt rR_{J(d)}.
 \end{aligned}
\]

By \(\texttt{lem:conventional-hac-consistency}\),
\[
 \max_{s\in\mathcal S}|\widehat V_s^{\mathrm{conv}}-V_s|
 \to_{P_n}0,
 \qquad
 \Pr_{P_n}\left\{\min_{s\in\mathcal S}
 \widehat V_s^{\mathrm{conv}}\geq\frac\nu2\right\}\to1.
\]
Selection cannot enlarge the first error, since
\[
 |\widehat V_{\widehat J}^{\mathrm{conv}}-V_{\widehat J}|
 \leq\max_{s\in\mathcal S}
 |\widehat V_s^{\mathrm{conv}}-V_s|.
\]
On the displayed positive-variance event, the inverse-square-root map is Lipschitz with a constant depending only on \(\nu\).  The selected numerator is tight, so replacing \(V_{\widehat J}^{-1/2}\) by \((\widehat V_{\widehat J}^{\mathrm{conv}})^{-1/2}\) costs \(o_{P_n}(1)\).  Reapplying the same null-boundary switching argument and using \(\widehat{\operatorname{se}}_s^{\mathrm{conv}}=\sqrt{\widehat V_s^{\mathrm{conv}}/n}\) gives
\[
 \frac{\widehat\tau_{\mathrm{sel}}-\tau}
 {\widehat{\operatorname{se}}^{\mathrm{conv}}_{\widehat J}}
 \Rightarrow
 \frac{A_1+\sqrt rR_1}{\sqrt{V_1}}\mathbf 1\{d+D\leq0\}
 +\frac{A_2+\sqrt rR_2}{\sqrt{V_2}}\mathbf 1\{d+D>0\}.
\]
Because \(X=D/\sigma_D\) and \(k=-d/\sigma_D\), the selection event \(d+D\leq0\) is exactly \(X\leq k\).  The last display is therefore
\[
 T_1\mathbf 1\{X\leq k\}+T_2\mathbf 1\{X>k\}.
\]

For the qualified reduction, first fix role \(s\) before observing the data.  The switching indicators disappear and the same linearization and variance consistency give \(T_s\sim N(0,1)\), since \(V_s=\operatorname{Var}(A_s+\sqrt rR_s)>0\).  On a separate sequence satisfying \(|\sqrt n(L_1-L_2)|\to\infty\), the centered loss fluctuation supplied by \(\texttt{lem:ridge-loss-expansion}\) is tight, so the sign of the population gap selects one role with probability tending to one and yields the same fixed-role limit.  Proposition \(\texttt{prop:shi-anchor-transfer}\) proves that the explicit stationary square family is a common subclass of this reduced-form setup and the published fixed-role setup.  Hence the reduction is valid on that intersection; no inclusion between the two full classes, and no statement about the rectangular application rule, follows.
%%% FIELD panel-class-replacement
For each fixed \(D_{\mathrm{gap}}<\infty\),
\[
 \mathcal M(D_{\mathrm{gap}})
 =\{(P_n)_{n\geq1}:(P_n)_{n\geq1}\text{ satisfies }
 \text{ass:stationarity},\text{ass:geometric-beta-mixing},
 \text{ass:moment-bound},\text{ass:sample-ratio},
 \text{ass:bridge-rank},\text{ass:bridge-centering},
 \text{ass:bounded-selector-gap-array},\text{ass:covariance-limit},
 \text{ass:att-variance-margin},\text{ and }
 \text{ass:selector-variance-margin}\}.
\]
This is the deliberately broad, unstabilized whole-array class: a uniform finite-row lower bound for the late and post scalar block variances is not part of its membership conditions.
%%% FIELD stabilized-class-construction
For fixed \(D_{\mathrm{gap}}<\infty\), define
\[
 \begin{aligned}
 \mathcal M_{\mathrm{stab}}(D_{\mathrm{gap}})
 :=\bigg\{P_\bullet=(P_j)_{j\geq1}\in\mathcal M(D_{\mathrm{gap}}):\ &
 \inf_{\substack{j\geq1:\,|\mathcal B_j|\geq1\\s\in\mathcal S}}
 \operatorname{Var}_{P_j}\!\left(
  |\mathcal B_j|^{-1/2}\sum_{t\in\mathcal B_j}\psi_{s,t}
 \right)\geq\nu,\\
 &\inf_{\substack{j\geq1\\s\in\mathcal S}}
 \operatorname{Var}_{P_j}\!\left(
  m_j^{-1/2}\sum_{t\in\mathcal C_j}v_{s,t}
 \right)\geq\nu
 \bigg\}.
 \end{aligned}
\]
Both added requirements are finite-row variance inequalities for scalar functions of the primitive panel and its population bridge.  They do not assume a Gaussian approximation, convergence rate, nonsingularity of a covariance matrix stacked across roles, or consistency of a variance estimator.
%%% FIELD uniform-gap-statement
For every fixed \(D_{\mathrm{gap}}<\infty\), the early-select, logarithm-squared-gap, late-fit estimator obeys
\[
 \sup_{P_\bullet\in\mathcal M_{\mathrm{stab}}(D_{\mathrm{gap}})}
 \mathbb E_{P_n}\!\left[
  \sup_{x\in\mathbb R}
  \left|
   \Pr_{P_n}\!\left\{
    \frac{\widehat\tau_{\mathrm{gap}}-\tau}
         {\widehat{\operatorname{se}}_{\widehat J_A}}
    \leq x\,\middle|\,\mathcal F_A
   \right\}-\Phi(x)
  \right|
 \right]\longrightarrow0.
\]
The conclusion remains valid when the covariance matrix obtained by stacking the two roles is singular.  The class is nonempty: if \(0<\kappa<2\), \(0<\nu\leq3/2\), and \(C_{\mathrm{mom}}\) exceeds the required Gaussian moment, then for every \(D_{\mathrm{gap}}\geq0\) the tied constant-treatment Gaussian array at
\[
 \theta_0=(1,1,1/2,1,1,1/2,r,1,1)
\]
in \(\texttt{def:primitive-gaussian-family}\) belongs to \(\mathcal M_{\mathrm{stab}}(D_{\mathrm{gap}})\) and to the published fixed-role common subclass certified by \(\texttt{prop:shi-anchor-transfer}\).
%%% FIELD uniform-gap-proof
Write
\[
 a_n=\lfloor\varrho n\rfloor,
 \qquad b_n=\lceil(\log n)^2\rceil,
 \qquad N_n=|\mathcal B_n|=n-a_n-b_n.
\]
Because \(0<\varrho<1\), one has \(N_n/n\to1-\varrho\in(0,1)\).  Assumption \(\texttt{ass:sample-ratio}\) gives \(n/m\to r\in(0,\infty)\), so \(N_n\), \(m\), and \(n\) are mutually comparable for all sufficiently large \(n\).  All constants below are uniform over \(P_\bullet\in\mathcal M_{\mathrm{stab}}(D_{\mathrm{gap}})\), over the two roles, and over sufficiently large \(n\).

We first obtain uniform moment and covariance bounds from primitive conditions.  Assumption \(\texttt{ass:bridge-rank}\) gives \(\lVert M_s^{-1}\rVert_{\mathrm{op}}\leq\kappa^{-1}\).  By H\"older's inequality and \(\texttt{ass:moment-bound}\),
\[
 \lVert\mathbb E_n(Z_{s,t}Y_t(0))\rVert_2
 +\lVert\mu_s\rVert_2\leq C.
\]
Consequently
\[
 \lVert\alpha_s\rVert_2+\lVert H_s\rVert_2\leq C.
\]
The score \(\psi_{s,t}=H_sZ_{s,t}\{Y_t(0)-W_{D_s,t}^{\prime}\alpha_s\}\) is quadratic in primitive coordinates with uniformly bounded \((4+\delta_0/2)\)-moment, while \(v_{s,t}\) has a uniformly bounded \((8+\delta_0)\)-moment.  Assumption \(\texttt{ass:bridge-centering}\), the definition of \(\alpha_s\), and the definition \(\tau=\mathbb E_n\Delta_t\) give
\[
 \mathbb E_n u_{s,t}=0,
 \qquad
 \mathbb E_n\psi_{s,t}=0,
 \qquad
 \mathbb E_n v_{s,t}=0.
\]
The last equality is the identifying calculation
\[
 \mathbb E_n v_{s,t}
 =\mathbb E_n\{Y_t(0)-W_{D_s,t}^{\prime}\alpha_s\}
  +\mathbb E_n\Delta_t-\tau=0;
\]
the observed post residual therefore targets the causal quantity only because the two declared bridge and treatment-effect equalities hold.

For completeness, let \(U\) and \(V\) be centered functions of two time sets separated by \(j\) observations, with common \((2+\epsilon)\)-moment bound.  Truncate each at \(R\).  The definition of the absolute-regularity coefficient bounds the covariance of the bounded parts by \(4R^2\beta_n(j)\).  H\"older's and Markov's inequalities bound the sum of the tail terms by \(CR^{-\epsilon}\).  Choosing \(R=\beta_n(j)^{-1/(4+2\epsilon)}\) yields, after a harmless weakening of the exponent,
\[
 |\operatorname{Cov}_{P_n}(U,V)|
 \leq C\beta_n(j)^{\epsilon/(4+2\epsilon)}
 \leq C_1e^{-c_1j}.
\]
Thus every scalar score used below has an absolutely summable autocovariance sequence, uniformly in the row law.  Define the row-specific long-run variances
\[
 \Omega_{\psi,s,n}=\sum_{j\in\mathbb Z}
 \operatorname{Cov}_{P_n}(\psi_{s,0},\psi_{s,j}),
 \qquad
 \Omega_{v,s,n}=\sum_{j\in\mathbb Z}
 \operatorname{Cov}_{P_n}(v_{s,0},v_{s,j}).
\]
For either stationary centered scalar process \(w_t\), with autocovariance \(\gamma_{w,n}(j)\), direct counting gives
\[
 \operatorname{Var}_{P_n}\!\left(N^{-1/2}\sum_{t=1}^Nw_t\right)
 =\gamma_{w,n}(0)+2\sum_{j=1}^{N-1}\left(1-\frac jN\right)\gamma_{w,n}(j).
\]
The exponential covariance envelope therefore implies
\[
 \left|
  \operatorname{Var}_{P_n}\!\left(N^{-1/2}\sum_{t=1}^Nw_t\right)
  -\sum_{j\in\mathbb Z}\gamma_{w,n}(j)
 \right|
 \leq \frac2N\sum_{j\geq1}j|\gamma_{w,n}(j)|
       +2\sum_{j\geq N}|\gamma_{w,n}(j)|
 \leq \frac CN.
\]
Applying this equality to the two variance inequalities in \(\texttt{def:panel-stabilized-class}\) shows, for all large \(n\),
\[
 \min_{s\in\mathcal S}\Omega_{\psi,s,n}\geq\frac\nu2,
 \qquad
 \min_{s\in\mathcal S}\Omega_{v,s,n}\geq\frac\nu2.
\]
The moment bounds give corresponding uniform upper bounds.

We next derive the late-fit linearization, including every product remainder.  Let bars with subscripts \(B\) and \(C\) denote averages over \(\mathcal B_n\) and \(\mathcal C_n\).  For a fixed role, the late empirical moment equation gives the exact identity
\[
 \widehat\alpha_{B,s}-\alpha_s
 =\widehat M_{B,s}^{-1}\overline u_{B,s}
\]
whenever \(\widehat M_{B,s}\) is invertible.  Exponential covariance summability and Chebyshev's inequality give
\[
 \widehat M_{B,s}-M_s=O_{P_n}(N_n^{-1/2}),
 \qquad
 \overline u_{B,s}=O_{P_n}(N_n^{-1/2})
\]
uniformly.  Hence
\[
 \sup_{P_\bullet}\Pr_{P_n}\!\left\{
  \lVert\widehat M_{B,s}-M_s\rVert_{\mathrm{op}}>\frac\kappa2
 \right\}\longrightarrow0.
\]
On the complementary event, \(\lVert\widehat M_{B,s}^{-1}\rVert_{\mathrm{op}}\leq2/\kappa\), and the resolvent identity yields
\[
 \widehat\alpha_{B,s}-\alpha_s
 =M_s^{-1}\overline u_{B,s}+O_{P_n}(N_n^{-1}).
\]
The fixed measurable fallback is confined to the vanishing-probability event just displayed.

The role-specific post residual average satisfies exactly
\[
 \widehat\tau_{\mathrm{gap},s}-\tau
 =\overline v_{C,s}
  -\overline W_{C,s}^{\prime}(\widehat\alpha_{B,s}-\alpha_s).
\]
Since \(\overline W_{C,s}-\mu_s=O_{P_n}(m^{-1/2})\), substituting the preceding bridge expansion and using \(H_s=-\mu_s^{\prime}M_s^{-1}\) gives
\[
 \widehat\tau_{\mathrm{gap},s}-\tau
 =\overline v_{C,s}+\overline\psi_{B,s}
  +O_{P_n}(N_n^{-1})
  +O_{P_n}(m^{-1/2}N_n^{-1/2}).
\]
The horizon bounds make the last two terms \(o_{P_n}(n^{-1/2})\), uniformly for both roles.  Set
\[
 G_{s,n}=\sqrt n(\overline\psi_{B,s}+\overline v_{C,s}),
 \qquad
 W_{s,n}=\frac n{N_n}\Omega_{\psi,s,n}
          +\frac nm\Omega_{v,s,n}.
\]
The preceding lower bounds and horizon comparability imply
\[
 0<c\leq\min_{s\in\mathcal S}W_{s,n}
 \leq\max_{s\in\mathcal S}W_{s,n}\leq C<\infty.
\]

We now prove the required scalar Gaussian approximation without inverting a covariance matrix across roles.  Put \(N_{*,n}=\min(N_n,m)\),
\[
 \ell_n=\lfloor N_{*,n}^{1/4}\rfloor,
 \qquad
 g_n=\lceil(\log N_{*,n})^2\rceil.
\]
For the proof only, partition the late and post observations into big blocks of length \(\ell_n\), separated by small blocks of length \(g_n\), and omit at most one incomplete big block at each endpoint and at the late--post boundary.  Each observation in \(G_{s,n}\) has deterministic weight bounded by \(C n^{-1/2}\).  Exponential covariance summability therefore bounds the variance of all omitted terms by
\[
 C\left(\frac{g_n}{\ell_n}+\frac{\ell_n}{n}\right)=o(1).
\]
If \(K_n=O(n/\ell_n)\) is the number of retained big blocks, telescoping the defining absolute-regularity bound over successive separated blocks shows that their joint characteristic function differs from the product of their marginal characteristic functions by at most
\[
 C K_n\beta_n(g_n)
 \leq C\frac n{\ell_n}\exp\{-c_\beta(\log N_{*,n})^2\}=o(1).
\]
No nonsingularity is used in this factorization.

Let \(S_{j,n}\) denote one centered weighted big-block sum.  The elementary inequality
\[
 \left|\sum_{t=1}^{\ell_n}x_t\right|^4
 \leq\ell_n^3\sum_{t=1}^{\ell_n}|x_t|^4
\]
and the uniform fourth-moment bound give
\[
 \sum_{j=1}^{K_n}\mathbb E_n|S_{j,n}|^4
 \leq C\frac n{\ell_n}\frac{\ell_n^4}{n^2}
 =C\frac{\ell_n^3}{n}=o(1).
\]
Also,
\[
 \sum_j\mathbb E_n|S_{j,n}|^3
 \leq
 \left(\sum_j\mathbb E_nS_{j,n}^2\right)^{1/2}
 \left(\sum_j\mathbb E_nS_{j,n}^4\right)^{1/2}
 =o(1).
\]
Taylor's formula for each marginal characteristic function, multiplication of the factors, and the last display therefore show that the retained independent-block comparison has characteristic function
\[
 \exp\left\{-\frac{u^2}{2}\sum_j\operatorname{Var}_{P_n}(S_{j,n})\right\}+o(1)
\]
at every fixed \(u\).  Restoring the omitted blocks changes the characteristic function by \(o(1)\).  Finally, the covariance between the adjacent late and post sums is bounded by
\[
 \left|\operatorname{Cov}_{P_n}\!\left(
  \sum_{t\in\mathcal B_n}\psi_{s,t},
  \sum_{t\in\mathcal C_n}v_{s,t}
 \right)\right|
 \leq C\sum_{j\geq1}j e^{-c_1j}\leq C,
\]
so after the weights in \(G_{s,n}\) are applied it is \(O(n^{-1})\).  The finite-block-to-long-run-variance bound already proved gives
\[
 \operatorname{Var}_{P_n}(G_{s,n})=W_{s,n}+o(1)
\]
uniformly.  Since \(W_{s,n}\) is bounded above and away from zero, the characteristic-function calculation proves
\[
 \sup_{P_\bullet\in\mathcal M_{\mathrm{stab}}(D_{\mathrm{gap}})}
 \sup_{x\in\mathbb R}
 \left|
  \Pr_{P_n}\{G_{s,n}/\sqrt{W_{s,n}}\leq x\}-\Phi(x)
 \right|\longrightarrow0.
\]
To justify the uniform conclusion explicitly, if it failed one could choose a violating sequence of row laws.  The preceding bounds continue to hold along that sequence and force convergence of every characteristic function to that of \(N(0,1)\); continuity of the normal distribution function then forces uniform convergence of the distribution functions, a contradiction.  This is coordinatewise in \(s\), so a singular covariance after stacking roles is irrelevant.

It remains to prove fitted-score Bartlett consistency rather than assume it.  For a true centered scalar score \(w_t\), let
\[
 \widetilde\gamma_{w,n}(j)=N^{-1}\sum_{t=j+1}^Nw_tw_{t-j}.
\]
For two lag-product terms whose leading indices differ by at most \(2j\), H\"older's inequality gives a constant covariance bound.  Beyond that range their underlying time sets are separated, and the truncation argument above gives an exponentially decreasing covariance bound.  Summation over the two indices yields
\[
 \operatorname{Var}_{P_n}\{\widetilde\gamma_{w,n}(j)\}
 \leq C\frac{j+1}{N}.
\]
Therefore the stochastic error of the Bartlett sum through bandwidth \(h\) is, in \(L^1\), at most
\[
 C\sum_{j=0}^h\sqrt{\frac{j+1}{N}}
 \leq C\frac{h^{3/2}}{\sqrt N}.
\]
The deterministic lag-truncation, Bartlett-weight, and finite-sample normalization errors are bounded by
\[
 C\sum_{j>h}e^{-c_1j}
 +\frac C{h+1}\sum_{j=1}^hj e^{-c_1j}
 +\frac CN\sum_{j=1}^hj e^{-c_1j}
 \leq C(e^{-c_1h}+h^{-1}+N^{-1}).
\]
Replacing the unknown zero mean by the sample mean costs \(O_{P_n}(hN^{-1/2})\), by the root-\(N\) mean bound and Cauchy--Schwarz.  With \(h=h_n=\lfloor N_{*,n}^{1/5}\rfloor\), every displayed error tends to zero uniformly.  Hence the oracle Bartlett estimators converge in probability to \(\Omega_{\psi,s,n}\) and \(\Omega_{v,s,n}\), uniformly over roles and laws.

For the fitted scores, the inverse event and expansions above imply
\[
 \widehat H_s-H_s=O_{P_n}(n^{-1/2}),
 \qquad
 \left\{N_n^{-1}\sum_{t\in\mathcal B_n}
  (\widehat\psi_{s,t}-\psi_{s,t})^2\right\}^{1/2}
 =O_{P_n}(n^{-1/2}),
\]
and
\[
 \left\{m^{-1}\sum_{t\in\mathcal C_n}
  (\widehat v_{s,t}-v_{s,t})^2\right\}^{1/2}
 =O_{P_n}(n^{-1/2}).
\]
These follow by expanding the fitted score differences into coefficient errors times \(u_{s,t}\), \(Z_{s,t}W_{D_s,t}^{\prime}\), and \(W_{D_s,t}\), whose empirical second moments are uniformly bounded in probability.  The centered Bartlett quadratic form has operator norm at most \(2h_n+1\); Cauchy--Schwarz therefore bounds the oracle-to-fitted replacement by \(O_{P_n}(h_nn^{-1/2})=o_{P_n}(1)\).  Consequently, jointly for the two roles,
\[
 \widehat\Omega_{\psi,s}-\Omega_{\psi,s,n}=o_{P_n}(1),
 \qquad
 \widehat\Omega_{v,s}-\Omega_{v,s,n}=o_{P_n}(1)
\]
uniformly.  The scalar lower bounds imply
\[
 \frac{
  \widehat{\operatorname{se}}_s^2
 }{
  \Omega_{\psi,s,n}/N_n+\Omega_{v,s,n}/m
 }\to_{P_n}1
\]
uniformly, and the probability of an empirical inverse or nonpositive-variance fallback tends uniformly to zero.  Combining this ratio, the late-fit linearization, and the scalar Gaussian approximation gives, for the fixed-role future statistic
\[
 K_{s,n}=\frac{\widehat\tau_{\mathrm{gap},s}-\tau}
 {\widehat{\operatorname{se}}_s},
\]
the bound
\[
 \epsilon_n:=
 \max_{s\in\mathcal S}
 \sup_{P_\bullet\in\mathcal M_{\mathrm{stab}}(D_{\mathrm{gap}})}
 \sup_{x\in\mathbb R}
 |\Pr_{P_n}\{K_{s,n}\leq x\}-\Phi(x)|\longrightarrow0.
\]

Finally, \(K_{s,n}\) is measurable with respect to the late and post observations.  Those observations are separated from \(\mathcal F_A\) by \(b_n\) discarded observations.  From the definition of absolute regularity,
\[
 \mathbb E_{P_n}\sup_{H}
 |\Pr_{P_n}(H\mid\mathcal F_A)-\Pr_{P_n}(H)|
 \leq2\beta_n(b_n)
 \leq2C_\beta e^{-c_\beta(\log n)^2},
\]
where the supremum is over events measurable in the late and post observations.  Since \(\widehat J_A\) is \(\mathcal F_A\)-measurable,
\[
 \Pr\{K_{\widehat J_A,n}\leq x\mid\mathcal F_A\}
 =\sum_{s=1}^2\mathbf 1\{\widehat J_A=s\}
   \Pr\{K_{s,n}\leq x\mid\mathcal F_A\}.
\]
Taking the supremum over \(x\), then expectation, and using the last two displays gives
\[
 \sup_{P_\bullet\in\mathcal M_{\mathrm{stab}}(D_{\mathrm{gap}})}
 \mathbb E_{P_n}\sup_x
 |\Pr\{K_{\widehat J_A,n}\leq x\mid\mathcal F_A\}-\Phi(x)|
 \leq4C_\beta e^{-c_\beta(\log n)^2}+\epsilon_n\longrightarrow0.
\]
No selection-probability denominator appears, so the conclusion covers exact and local selector ties.

To prove nonemptiness, use the tied constant-treatment Gaussian point \(\theta_0\) from \(\texttt{def:primitive-gaussian-family}\).  At that point \(B=2\), \(K=1/2\), the bridge matrices equal \(M_s=2\), and the calculations in the primitive family give \(\operatorname{Var}(D)=13/4\).  Write \(e_0^\star=e_0+\xi\), so \(\operatorname{Var}(e_0^\star)=1\).  For either role,
\[
 \psi_{s,t}=-\frac12(F_t+e_{3-s,t})(e_0^\star-e_{s,t}),
 \qquad
 v_{s,t}=e_0^\star-e_{s,t}.
\]
The two factors in the expression for \(\psi_{s,t}\) are independent and centered in product because the second factor is centered.  Therefore independence over time gives, for every nonempty block,
\[
 \operatorname{Var}(\psi_{s,t})
 =\frac14\mathbb E(F_t+e_{3-s,t})^2
       \mathbb E(e_0^\star-e_{s,t})^2
 =\frac14(3)(2)=\frac32,
 \qquad
 \operatorname{Var}(v_{s,t})=2.
\]
Thus both finite-row stabilization inequalities hold whenever \(0<\nu\leq3/2\); the rank margin holds for \(0<\kappa<2\), and the exact loss tie belongs to every envelope \(D_{\mathrm{gap}}\geq0\).  Gaussian moments meet the declared moment bound when \(C_{\mathrm{mom}}\) is chosen as in the statement.  Proposition \(\texttt{prop:shi-anchor-transfer}\) verifies the published fixed-role assumptions for this same constant-treatment array.  Hence the stabilized class and its common-subclass interpretation are nonvacuous.
